\section{The Ergodic Theorem/Law of Large Numbers}

\begin{theorem}[The Ergodic Theorem for Markov Chains]
    Let \(X \sim \Markov\pqty{\lambda, P}\) be irreducible, aperiodic and positive recurrent.

    Then, for any \(x \in I\), we have
    \[
        \frac{1}{n} \sum_{i = 0}^{n - 1} \indicator \set{X_i = x} \toAlmostSure \pi\pqty{x}.
    \]
\end{theorem}

Let \(\functiondomain{f}{I}{\RR}\) be any function such that
\[
    \Expt_{\pi} \bqty{\abs{f\pqty{X_0}}} < \infty,
\]
i.e.,
\[
    \sum_{x \in I} \pi\pqty{x} \abs{f\pqty{x}} < \infty,
\]
then
\begin{align*}
    \frac{1}{n} \sum_{i = 0}^{n - 1} f\pqty{X_i} &= \frac{1}{n} \sum_{i = 0}^{n - 1} \sum_{x \in I} f\pqty{x} \indicator \set{X_i = x} \\
    &= \sum_{x \in I} \bqty{f\pqty{x} \frac{1}{n} \sum_{i = 0}^{n - 1} \indicator \set{X_i = x}} \\
    &\toAlmostSure \sum_{x \in I} \bqty{f\pqty{x} \pi\pqty{x}} \\
    &= \Expt_{\pi} \bqty{f\pqty{X_0}}.
\end{align*}

\small

\begin{remark}[Non-Examinable]
    What the theorem says is that
    \[
        \Prob\pqty{\frac{1}{n} \sum_{i = 0}^{n - 1} \indicator \set{X_i = x} \to \pi\pqty{x}} = 1,
    \]
    for every \(x \in I\). It is a simple measure theoretic fact (see \textit{II Probability and Measure}) that since \(I\) is countable, we have
    \[
        \Prob\pqty{\forall x \in I: \frac{1}{n} \sum_{i = 0}^{n - 1} \indicator \set{X_i = x} \to \pi\pqty{x}} = 1.
    \]
\end{remark}

\begin{proof}[\proofname\ (Non-Examinable)]
    Recall the following notation. Fix \(x\), and let
    \[
        V_n\pqty{x} = \sum_{i = 0}^{n - 1} \indicator \set{X_i = x},
    \]
    and consider the recurrence times \(T_x^{\pqty{0}} = 0\), and
    \[
        T_x^{\pqty{n + 1}} = \inf \set{n > T_x^{\pqty{n}}}{X_n = x}
    \]
    if the chain starts at \(X_0 = x\). What we want to show is that
    \[
        \frac{1}{n} V_n\pqty{x} \toAlmostSure \pi\pqty{x}.
    \]

    But we know that \(\Prob\pqty{T_x < \infty} = 1\), so by the Strong Markov Property,
    \[
        \pqty{X_{T_n + k}}_{k \geq 0} \sim \Markov\pqty{\kronecker_x, P},
    \]
    and is independent if \(\pqty{X_0, \ldots, X_{T_n}}\).

    Therefore, the long-term behaviour of \(\flatfrac{V_n\pqty{x}}{n}\) is independent of the initial distribution \(\lambda\). So without loss of generality, we can assume that \(\lambda = \kronecker_x\).

    Define the \term{successive excursion length} \(S_x^{\pqty{k}}\) for \(k \geq 1\), by
    \[
        S_x^{\pqty{k}} = \begin{cases}
            T_x^{\pqty{k}} - T_x^{\pqty{k - 1}}, & \qif T_x^{\pqty{k - 1}} < \infty, \\
            \infty, & \quad \text{otherwise}.
        \end{cases}
    \]

    Then, by the definition, we have
    \[
        T_x^{\pqty{V_n\pqty{x} - 1}} \leq n - 1,
    \]
    or equivalently, by telescoping,
    \[
        S_x^{\pqty{1}} + \cdots + S_x^{\pqty{V_n\pqty{x} - 1}} \leq n - 1.
    \]

    Similarly, we have
    \[
        T_x^{\pqty{V_{n} \pqty{x}}} \geq n,
    \]
    and hence by telescoping,
    \[
        S_x^{\pqty{1}} + \cdots + S_x^{\pqty{V_n\pqty{x}}} \geq n.
    \]

    Hence,
    \[
        S_x^{\pqty{1}} + \cdots + S_x^{\pqty{V_n\pqty{x} - 1}} \leq n \leq S_x^{\pqty{1}} + \cdots + S_x^{\pqty{V_n\pqty{x}}}.
    \]

    Then, we have
    \[
        \frac{1}{V_n\pqty{x}} \sum_{j = 1}^{V_n\pqty{x} - 1} S_x^{\pqty{j}} \leq \frac{n}{V_n\pqty{x}} \leq \frac{1}{V_n\pqty{x}} \sum_{j = 1}^{V_n\pqty{x}} S_x^{\pqty{j}}.
    \]

    We have \(V_n\pqty{x} \toAlmostSure \infty\) as \(n \to \infty\), and the process \(\pqty{S_x^{\pqty{k}}}_{k \geq 1}\) is i.i.d. with mean
    \[
        \Expt_x\pqty{S_x^{\pqty{n}}} = \Expt_x\pqty{T_x} = \frac{1}{\pi\pqty{x}}.
    \]

    Hence, by the Strong Law of Large Numbers (see \textit{IA Probability}), both sides of the inequality converge almost surely to \(\Expt_x\pqty{S_x^{\pqty{n}}} = \flatfrac{1}{\pi\pqty{x}}\). Hence, we have
    \[
        \frac{n}{V_n\pqty{x}} \toAlmostSure \frac{1}{\pi\pqty{x}},
    \]
    which is what we want to prove.
\end{proof}